\clearpage
\SourceLeaf{97}%

\begin{center}
\large [\"UBER DAS WESEN UND DIE DEFINITION DER FUNCTIONEN.]
\end{center}

\begin{center}
\textsc{Gauss} an \textsc{Bessel}, 18. December 1811.
\end{center}

\noindent\hspace*{\parindent}$\mathord{.}\,\mathord{.}\,\mathord{.}\,\mathord{.}\,\mathord{.}\,\mathord{.}$ Zuv\"orderst w\"urde ich jemand, der eine neue Function in die Analyse einf\"uhren will, um eine Erkl\"arung bitten, ob er sie schlechterdings bloss auf reelle Gr\"ossen (reelle Werthe des Arguments der Function) angewandt wissen will, und die imagin\"aren Werthe des Arguments gleichsam nur als ein \"Uberbein ansieht --- oder ob er meinem Grundsatze beitrete, dass man in dem Reiche der Gr\"ossen die imagin\"aren $a+b\sqrt{-1}=a+bi$ als gleiche Rechte mit den reellen geniessend ansehen m\"usse. Es ist hier nicht von practischem Nutzen die Rede, sondern die Analyse ist mir eine selbst\"andige Wissenschaft, die durch Zur\"ucksetzung jener fingirten Gr\"ossen ausserordentlich an Sch\"onheit und Rundung verlieren und alle Augenblick Wahrheiten, die sonst allgemein gelten, h\"ochst l\"astige Beschr\"ankungen beizuf\"ugen gen\"othigt sein w\"urde. $\mathord{.}\,\mathord{.}\,\mathord{.}\,\mathord{.}\,\mathord{.}\,\mathord{.}\,\mathord{.}\,\mathord{.}\,\mathord{.}$ Was soll man sich nun bei $\int \varphi x\mathord{.} dx$ f\"ur $x=a+bi$ denken? Offenbar, wenn man von klaren Begriffen ausgehen will, muss man annehmen, dass $x$ durch unendlich kleine Incremente (jedes von der Form $a+\beta i$) von demjenigen Werthe, f\"ur welchen das Integral 0 sein soll, bis zu $x=a+bi$ \"ubergeht und dann alle $\varphi x\mathord{.} dx$ \textls{summirt}. So ist der Sinn vollkommen festgesetzt. Nun aber kann der \"Ubergang auf unendlich viele Arten geschehen: so wie man sich das \textls{ganze Reich aller reellen Gr\"ossen} durch eine unendliche gerade Linie denken kann, so kann man das \textls{ganze Reich aller} Gr\"ossen, reeller und
